\documentclass[../main]{subfiles}
\begin{document}
\ifSubfilesClassLoaded{\setcounter{page}{40}}{}
\section{Communism and the Management of Social Life}
\label{sec:07_communism_management}

The production and distribution of free time during the period of transition to communism require the state. However, the state is necessary only insofar as the production of free time coincides with the production and distribution of the things people need, and insofar as this production of things is not yet carried out in a communist way. Free time itself, emerging as genuinely free (the time of communism), escapes any state management external to it.

Within the space of free time, the state does not exist even now, under capitalism. Wherever free communist activity takes place, management is carried out by all the members of a free association of people. As the share of communist activity grows, the state withers away.

Promoting the production of free time and of directly collective forms of free activity is an act of abolishing the conditions for the existence of classes during the period of transition, and to the same extent an act of the self-abolition of the state. The state is merely an apparatus for securing the interests of the ruling class, exercising management in the interests of that class, and this primary function falls away as classes, including the ruling class, are abolished.

The task of the state during the period of transition thus comes down to ensuring the restriction (abolition) of the commodity character of production and to concern for creating the conditions for the free communist activity of the members of society.

The first of these tasks confronts capitalist states as well. Capitalism constantly breeds wars, and the logic of war demands a system of direct distribution of products intended for personal consumption. A state that fails to cope with this task is doomed to destruction (either on its own, as a state, or together with its population). The task of establishing such a system entails, at a number of points, the curbing of individual capitalists in the interests of capitalists as a class, which cannot always be achieved in full or with a sufficient degree of effectiveness.

In war, personal consumption—of both food and clothing supplies—cannot be provided for through the capitalist private character of appropriation of the products of labour. A soldier, even a mercenary, must be supplied by his employer during wartime with, at the very least, weapons, food and uniforms. Otherwise, the army degenerates into gangs of brigands. But such gangs are incapable of waging modern warfare. To wage war, a modern army needs high-tech weaponry, the means of personal consumption necessary for reproducing the living personality of the soldier, and personal protective equipment. And even if all of this is secured through the outright plunder of the population of captured territories, it cannot be the personal affair of individual soldiers. It is merely a condition that enables the soldier to do his job. But this condition differs from the conditions of wage labour, which is based on the private character of appropriation. In classical capitalist relations of the purchase and sale of labour power, the appropriation of the products of labour, in all of its aspects that are essential and significant for the social reproduction of the individual, is placed outside the sphere of production based on wage labour. Personal consumption is determined by the private character of appropriation of the products of labour, mediated by money. And this, in turn, guarantees that the worker will be compelled, again and again, to return to work in order to provide himself with everything he needs. In war, things are radically different. Here, consumption cannot be left to individuals themselves, even formally.

Thus, modern warfare is, among other things, a competition between forms of non-commodity production that took shape under capitalism itself—the direct production of man. This applies both to the army as a military force proper and to the labour army, which is formed in various ways. The emphasis shifts. Whereas in the normal course of capitalism what comes to the fore is that a person must do his work in order to live, here, on the contrary, a person must live in order to do his work. The state is forced to take responsibility for this. Moreover, commodity-money relations become an obvious obstacle to the performance of these functions.

War makes the old, pre-war ways and forms of human activity impossible. Under such conditions, the state objectively grows stronger, takes on the functions of distribution in the process of the direct production of man, both in the army and on the home front, and works to establish and maintain economic ties. But the state, as an apparatus of suppression of one class by another—or, more precisely, of the majority by a minority—is in principle incapable of fully accomplishing these tasks, let alone setting itself tasks that do not serve the interests of the ruling class. Yet it is precisely tasks of this kind (those at odds with the interests of the ruling class) that multiply for society during war. They objectively confront people because their very physical existence depends on solving them. The state often does not even attempt to solve these problems, leaving them to a wide variety of non-state organisations and associations of people. This compels millions of people to seek new forms of organisation that would take over state functions. And these bodies of social self-government, which take on the functions of the state, strictly speaking, in their tendency both are and are not a state.

The most important and formative of the problems whose solution requires new organisational forms in opposition to the capitalist state is the prevention and ending of war, along with minimising the losses and destruction it inflicts on society. Accomplishing this crucial task runs directly counter to the interests of capital, which has not yet extracted from the war all the benefits it could. Let us stress that what may carry historical meaning here is not war as such, nor violence as such, but resistance to the violence of the reactionary class, violence carried out by factions of that class in the name of their own selfish interests, and by the hands of another class at that. It is in resistance and in the forms of resistance (though not only there), rather than in war as such, that the most important tendency towards self-organisation on the approaches to revolution lies. Without such self-organisation, the methods of the direct, unmediated by things and money, production of man that have been developed within capitalism, along with the economic mechanisms for implementing them, cannot become the starting point for society’s transition to new foundations of development.

These forms of self-organisation for managing public affairs, generated by the tendency towards war and by war itself (as opposition to war), are transitional and can move in three directions: 1) integration with the state; 2) degeneration and disappearance (integration does not always mean degeneration); 3) taking over all the functions of the state and turning into a new state power. These organisational forms are thus, objectively, also forms of class struggle wherever they are forced to oppose the state by taking over the functions of particular state bodies, and hence stripping the apparatus of the ruling class’s political domination of its powers. And while under certain conditions this actually comes as a relief to those bodies, since they themselves are unable to perform their functions, any further movement along this path inevitably runs up against the opposition of the state and its attempts to restrict the functions and influence of the population’s forms of self-organisation. The outcome and character of this confrontation will depend on the specific circumstances of the struggle and the actions of the contending parties.

Looking at the history of the twentieth century, one of the most effective forms of self-organisation to take on the functions of the state were the Soviets. The Soviets of Workers’, Soldiers’ and Peasants’ Deputies served simultaneously as representative-legislative, executive and supervisory bodies. And although the Soviets first appeared as early as 1905, their real development into a state force, dual power, let alone the slogan “All power to the Soviets” and the concrete steps from that slogan to its implementation, could only arise amid the political bankruptcy of the old state in the course of an exhausting war. After the February Revolution, the bankruptcy of the state as an apparatus claiming to organise the life of society became obvious, since not only the tsarist government but also the Provisional Government proved powerless in this respect. At the same time, the Soviets were rising and actively involving themselves in the affairs of both local self-government and state administration. A situation of dual power was taking shape both on the home front and in the army. For this reason, it is especially important to bear in mind that the situation of spring and early summer 1917 could arise again in the future. At that time, a peaceful transfer of power to the Soviets was possible as a step along the path towards social revolution, that is, towards creating the conditions for socialism. This possibility existed thanks to two circumstances: the strength and social significance of the Soviets themselves, and the fact that the population had been armed as a result of the war. The slogan “All power to the Soviets” was, at the time, a statement that such a peaceful path of revolution existed. Only after the July Days of 1917 did the peaceful path become impossible. But the very fact that such a possibility existed is what matters most to us. After all, if such conditions arise, it is vital not to miss them, especially since, given the development of modern weaponry, the peaceful path may well prove to be the only possible one for revolution.

Under present-day conditions, other forms can emerge, and are emerging, which are by no means always connected with war but are always connected with the need to solve important social problems. The essence of these transitional forms is the same everywhere: the population directly takes upon itself the functions of state power, thereby bringing about the transcendence of that power in the sense customary under capitalism.

Let us now turn to the management of social life under communism. In general, the management of social life comes down to planning and adjusting the plan in light of real changes, the expression of will (issuing directives), accounting and control over the managed process, and direct intervention where necessary.

Planning under communism is planning for the development of the human being, who becomes the principal aim of social production; it is planning for the production of new relations between people. Such planning coincides with forecasting, in the sense that the most reliable forecast is the one towards whose fulfilment systematic effort is directed. Society must learn to perform all these functions without an organ of “management” standing above it, which is in reality an organ of oppression, and which is called the state.

In the course of the becoming of the new collectivity, however, all the functions of the state that relate to the production of things must be handed over to things. In the period of transition to communism, the state does not simply wither away; it withers away only if it works towards its own withering away, to the extent that it takes care that people can do without it. All the remnants of unfreedom, including the ways of utilising time with which the existence of the state is bound up, stand opposed to free time as the space of human development, and stand opposed to free human activity. Where complete freedom is achievable, there is no place for the state. But where complete freedom is not yet directly possible, the existence of a certain type of state not only registers the fact of this impossibility but also serves as an instrument for overcoming it.

The same argument is advanced against the withering away of the state as against the sublation of the division of labour. Its essence is that social life is becoming very complex and is growing ever more complex, and that managing such a complex life therefore requires a differentiation of functions and a special apparatus. But it by no means follows from this complexity that special detachments of officials and armed men, set above society, are necessary. Were it not for the split of society into classes, a self-acting and self-governing association of people would be possible, for all its complexity corresponding to the level of complexity of social life, as Lenin, following Marx and Engels, insisted.

All the more so since, when we speak of communism, we mean a society in which the division of labour has been overcome, in which people are more or less equally developed in terms of health, sensuousness, thinking and morality, and have a polytechnical, scientific-theoretical, universal education at the level necessary for society today (that is, the highest), and can therefore master any particular “profession” in a short time and apply their efforts wherever they are needed and wherever they themselves find it interesting, even if one person devotes most of his life to studying birds and another to ballet.

But it is precisely for this reason that under communism everyone takes part in the management of social development, using the technologies achieved and relying on the culture of thinking that humanity has developed. This is the most important condition for the universal development of all members of society: throughout their lives, everyone must be able to manage, and be in the habit of managing, social affairs—all the affairs in which they themselves take part. Such a society is more realistic than the fantastical class peace and “harmonious relations” between classes propagandised by today’s bourgeois ideologists. On the contrary, Marxism insists that the very existence of the state shows that such peace is impossible as long as classes exist.

The withering away of the state requires the appropriate social and technical preconditions. As for the technical preconditions, with good universal polytechnical education grounded in modern technology, organising the management of social processes would not be all that difficult in a society without private property (which is why the formal socialisation of the means of production is the most important precondition). Of course, social barriers must be removed for this, which likewise demands serious effort from society and cannot be achieved overnight. The main point is that, purely technically, it no longer seems impossible even at the present stage. All information on the state of production and distribution in every sphere of social life in real time, with the possibility of tracing chains back to their very beginning, accounting of resources, needs and technological capabilities, the possibilities of communication and voting, and so on and so forth, from the lowest level (direct intervention) up to the highest level of collective decision-making, both in local communities and in society as a whole—all of it can be available in real time to every member of the association.

Voting and the other familiar expressions of democracy are needed only at first. With the transition from the government of persons to the administration of things and processes, the need for voting gradually falls away, giving way to the objective necessity inherent in things and in the interactions between them. Lenin pointed this out in The State and Revolution in his detailed examination of the legacy of Marx and Engels: “In speaking of the ‘withering away’, and the even more graphic and colourful ‘dying down of itself’, Engels refers quite clearly and definitely to the period after ‘the taking possession of the means of production in the name of society’, that is, after the socialist revolution. We all know that the political form of the ‘state’ at that time is the most complete democracy. But it never enters the head of any of the opportunists, who shamelessly distort Marxism, that Engels is consequently speaking here of democracy ‘dying down of itself’, or ‘withering away’. This seems very strange at first sight. But it is ‘incomprehensible’ only to those who have not thought about democracy also being a state and, consequently, also disappearing when the state disappears. Revolution alone can ‘abolish’ the bourgeois state. The state in general, i.e., the most complete democracy, can only ‘wither away’.”

In today’s era of rapid technological development, organising the withering away of the state does not seem such a difficult matter from a purely technical standpoint. Under such conditions, accounting and control in the administration of things and processes become a matter of technology, and their results become accessible to everyone. Here the need for hierarchy falls away, and the precondition for these functions to become separated into a distinct profession disappears. With the introduction of automated management systems, and with full access to databases and the ability to work with them guaranteed for everyone, the need disappears for a “special army of officials” who could stand, even “in appearance”, “above” society as a special “apparatus”.

In this way, one of the most important tasks of the transition to socialism will be accomplished: establishing accounting and control, the importance of which Lenin constantly stressed: “Accounting and control—that is mainly what is needed for the ‘smooth working’, for the proper functioning, of the first phase of communist society.” What we are talking about is accounting and control over social production of a kind whose very organisation already means the organisation of new forms of collectivity: “If really all take part in the administration of the state, capitalism cannot retain its hold. The development of capitalism, in turn, creates the preconditions that enable really ‘all’ to take part in the administration of the state. Some of these preconditions are: universal literacy, which has already been achieved in a number of the most advanced capitalist countries, then the ‘training and disciplining’ of millions of workers by the huge, complex, socialised apparatus of the postal service, railways, big factories, large-scale commerce, banking, etc., etc.”

We, for our part, can emphasise once again that these preconditions now also include the development of computer technology and the potential it holds, which in principle cannot be realised under capitalism, where there is no single social economy and the connection between producers is mediated by the market. This potential for organising the management of social life was explored in the work of Victor Mikhailovich Glushkov, who was developing the nationwide system of economic management in the USSR. Glushkov considered the management of social production as a whole by means of control systems in which everyone has access to a highly detailed, real-time picture of production, while wherever human decisions are not required, the system acts automatically according to schemes set in advance by people in line with social goals.

The more technology develops, the easier it will be to organise management and to ensure that all members of society take part in it. As for the substantive side, it requires a corresponding development of the members of society. The difficulty is that the movement towards communism has to begin not with the people of the communist tomorrow, but with people shaped under the conditions of the division of labour, who overcome their limitations through education and collective activity.

Polytechnical education becomes a necessity for the members of a free association. The idea of polytechnism has been worked out quite thoroughly, so we can safely refer the reader to the relevant literature. Here we will only note that in order to manage social processes, members of society must possess an integral scientific picture of the world and a universal ability to apply their knowledge. A modern official has absolutely no need for a scientific picture of the world, even in rough outline, to perform his function, and so it simply disappears of its own accord as something superfluous, even if he once had one. An official is merely a living, partial component of an obsolete state machine. “The society that will organise production anew on the basis of a free and equal association of the producers will put the whole state machine where it will then belong: into the museum of antiquities, by the side of the spinning wheel and the bronze axe,” said Engels.

Of course, the very task of managing social processes in such an association, and of consciously building social life, will demand a great deal of people, including well-developed dialectical thinking, well-developed sensuousness, morality, and an adequate level of literacy. But these are not unattainable heights accessible only to a chosen few; they are something that, under the right conditions, will become the possession of everyone. As people are drawn into active social life as its subjects, they will come to the forefront of human culture. This is how we understand Lenin’s famous aphorism that every cook must learn to govern the state. The only obstacle to this is the mode of activity of people today, which turns one person into a state official and another into a cook.

Communist society is a free association of people governing its own development according to a rational plan.

How exactly will it govern itself, and by what specific means and organisational measures? We cannot answer this question. That is a matter for the association itself, which will of necessity have to develop the appropriate organisational forms under concrete historical conditions. Self-government will develop along with the development of communist society.
\end{document}
